
\section{Advanced Activities}

The following activities are intended for strong master's or PhD students.
They emphasize the mathematical and computational differences among
single-pass sequential, iterative sequential, nested, and simultaneous
control co-design architectures.

\subsection*{Problem 1: Exact Comparison of Four CCD Architectures}

Consider the unconstrained control co-design problem
\begin{equation}
\min_{p,c}\;
J(p,c)
=
\frac{1}{2}
\left(
6p^2+6pc+4c^2
\right)
-8p-6c,
\end{equation}
where $p$ is a plant-design variable and $c$ is a controller-design
variable.

\begin{enumerate}
    \item Verify that $J$ is strictly convex.

    \item In a single-pass sequential design, first optimize $p$ using the
    nominal controller $\hat{c}=0$, and then optimize $c$ for the fixed
    value of $p$. Compute
    \begin{equation}
    p_{\mathrm{seq}},
    \qquad
    c_{\mathrm{seq}},
    \qquad
    J_{\mathrm{seq}}.
    \end{equation}

    \item Formulate the nested problem
    \begin{equation}
    \phi(p)=\min_c J(p,c),
    \qquad
    \min_p \phi(p).
    \end{equation}
    Derive $c^*(p)$ and the reduced objective $\phi(p)$.

    \item Solve the nested problem analytically.

    \item Solve the simultaneous problem by imposing
    \begin{equation}
    \nabla J(p,c)=\mathbf{0}.
    \end{equation}

    \item Show that the nested and simultaneous solutions are identical.

    \item Compute the relative performance loss of the single-pass
    sequential design:
    \begin{equation}
    \eta_{\mathrm{seq}}
    =
    \frac{J_{\mathrm{seq}}-J_{\mathrm{sim}}}
    {|J_{\mathrm{sim}}|}.
    \end{equation}

    \item Explain why strict convexity is important for the equivalence
    observed in this problem.
\end{enumerate}


\subsection*{Problem 2: Convergence of Iterative Sequential Design}

Consider the quadratic CCD problem
\begin{equation}
J(\mathbf{p},\mathbf{c})
=
\frac{1}{2}\mathbf{p}^TA\mathbf{p}
+
\mathbf{p}^TB\mathbf{c}
+
\frac{1}{2}\mathbf{c}^TC\mathbf{c}
-
\mathbf{a}^T\mathbf{p}
-
\mathbf{b}^T\mathbf{c},
\end{equation}
where $A$ and $C$ are symmetric positive-definite matrices.

An iterative sequential method performs
\begin{align}
\mathbf{p}^{(k+1)}
&=
\arg\min_{\mathbf{p}}
J\left(\mathbf{p},\mathbf{c}^{(k)}\right),
\\
\mathbf{c}^{(k+1)}
&=
\arg\min_{\mathbf{c}}
J\left(\mathbf{p}^{(k+1)},\mathbf{c}\right).
\end{align}

\begin{enumerate}
    \item Derive the explicit update equations
    \begin{align}
    \mathbf{p}^{(k+1)}
    &=
    A^{-1}
    \left(
    \mathbf{a}-B\mathbf{c}^{(k)}
    \right),
    \\
    \mathbf{c}^{(k+1)}
    &=
    C^{-1}
    \left(
    \mathbf{b}-B^T\mathbf{p}^{(k+1)}
    \right).
    \end{align}

    \item Let $\mathbf{c}^*$ denote the simultaneous optimum. Show that the
    controller error satisfies
    \begin{equation}
    \mathbf{c}^{(k+1)}-\mathbf{c}^*
    =
    M
    \left(
    \mathbf{c}^{(k)}-\mathbf{c}^*
    \right),
    \end{equation}
    where
    \begin{equation}
    M=C^{-1}B^TA^{-1}B.
    \end{equation}

    \item Prove that the iterative sequential method converges for every
    initial controller if
    \begin{equation}
    \rho(M)<1,
    \end{equation}
    where $\rho(\cdot)$ denotes spectral radius.

    \item Use
    \begin{equation}
    A=
    \begin{bmatrix}
    4&1\\
    1&3
    \end{bmatrix},
    \qquad
    C=
    \begin{bmatrix}
    3&0.5\\
    0.5&2
    \end{bmatrix},
    \end{equation}
    \begin{equation}
    B=
    \begin{bmatrix}
    1.8&0.4\\
    0.2&1.2
    \end{bmatrix},
    \qquad
    \mathbf{a}=
    \begin{bmatrix}
    2\\1
    \end{bmatrix},
    \qquad
    \mathbf{b}=
    \begin{bmatrix}
    1\\2
    \end{bmatrix}.
    \end{equation}

    Compute $\rho(M)$ and determine whether the iterative sequential
    method converges.

    \item Starting from
    \begin{equation}
    \mathbf{c}^{(0)}=
    \begin{bmatrix}
    0\\0
    \end{bmatrix},
    \end{equation}
    perform ten alternating iterations and compare the result with the
    simultaneous optimum.

    \item Explain how the strength of plant--controller coupling is reflected
    in $\rho(M)$.
\end{enumerate}


\subsection*{Problem 3: Nested LQR Control Co-Design}

Consider the scalar plant
\begin{equation}
\dot{x}=-px+u,
\qquad
x(0)=1,
\end{equation}
where
\begin{equation}
0.1\leq p\leq3
\end{equation}
is a plant-design variable.

Use the feedback controller
\begin{equation}
u=-Kx,
\qquad
K\geq0.
\end{equation}

The infinite-horizon system-level objective is
\begin{equation}
J(p,K)
=
\int_0^\infty
\left(
x(t)^2+r\,u(t)^2
\right)dt
+
\alpha p^2,
\end{equation}
with
\begin{equation}
r=0.2,
\qquad
\alpha=0.08.
\end{equation}

\begin{enumerate}
    \item Show that the closed-loop state is
    \begin{equation}
    x(t)=e^{-(p+K)t}.
    \end{equation}

    \item Derive the finite-dimensional objective
    \begin{equation}
    J(p,K)
    =
    \frac{1+rK^2}{2(p+K)}
    +
    \alpha p^2.
    \end{equation}

    \item For fixed $p$, solve the inner controller problem and show that
    \begin{equation}
    K^*(p)
    =
    -p+\sqrt{p^2+\frac{1}{r}}.
    \end{equation}

    \item Derive the nested value function
    \begin{equation}
    \phi(p)=J\left(p,K^*(p)\right).
    \end{equation}

    \item Independently derive the same inner solution using the scalar
    algebraic Riccati equation.

    \item Minimize $\phi(p)$ over the admissible interval and compute
    \begin{equation}
    p^*_{\mathrm{nested}},
    \qquad
    K^*_{\mathrm{nested}}.
    \end{equation}

    \item Solve the simultaneous problem
    \begin{equation}
    \min_{p,K}J(p,K)
    \end{equation}
    directly and verify that it produces the same optimum.

    \item Perform a single-pass sequential design by first setting $K=0$
    and optimizing $p$, and then optimizing $K$ for the fixed plant.
    Compare the sequential and coordinated solutions.
\end{enumerate}


\subsection*{Problem 4: Differentiating a Nested Value Function}

Consider the nested CCD problem
\begin{equation}
\min_{\mathbf{p}}\;
\phi(\mathbf{p})
=
\frac{1}{2}\mathbf{p}^TQ\mathbf{p}
+
\min_{\mathbf{c}}
\left[
\frac{1}{2}\mathbf{c}^TH\mathbf{c}
+
\mathbf{p}^TD\mathbf{c}
\right],
\end{equation}
subject to the inner equality constraint
\begin{equation}
A\mathbf{c}
=
\mathbf{b}+E\mathbf{p},
\end{equation}
where $H$ is symmetric positive definite and $A$ has full row rank.

\begin{enumerate}
    \item Write the KKT system for the inner problem:
    \begin{equation}
    \begin{bmatrix}
    H&A^T\\
    A&0
    \end{bmatrix}
    \begin{bmatrix}
    \mathbf{c}^*\\
    \boldsymbol{\lambda}^*
    \end{bmatrix}
    =
    \begin{bmatrix}
    -D^T\mathbf{p}\\
    \mathbf{b}+E\mathbf{p}
    \end{bmatrix}.
    \end{equation}

    \item Differentiate the KKT system with respect to $\mathbf{p}$ and
    derive
    \begin{equation}
    \frac{d\mathbf{c}^*}{d\mathbf{p}},
    \qquad
    \frac{d\boldsymbol{\lambda}^*}{d\mathbf{p}}.
    \end{equation}

    \item Derive the total derivative of the nested objective,
    \begin{equation}
    \frac{d\phi}{d\mathbf{p}}.
    \end{equation}

    \item Show how the envelope theorem eliminates the explicit
    $d\mathbf{c}^*/d\mathbf{p}$ term when the inner KKT conditions are
    satisfied exactly.

    \item Explain why this simplification may fail numerically when the inner
    problem is terminated with a large optimality or feasibility residual.

    \item Implement the nested gradient for a randomly generated
    positive-definite $H$ and verify it using complex-step or central finite
    differences.

    \item Investigate how the outer-gradient error changes when the inner
    solver tolerance is varied from $10^{-3}$ to $10^{-10}$.
\end{enumerate}


\subsection*{Problem 5: KKT Reformulation of a Constrained Bilevel CCD Problem}

Consider the nested problem
\begin{equation}
\min_{0\leq p\leq3}\;
\phi(p)
=
(p-2)^2
+
\min_{0\leq c\leq1}
\left[
(c-p)^2+0.1c^2
\right].
\end{equation}

\begin{enumerate}
    \item Solve the unconstrained inner problem and obtain
    \begin{equation}
    c_{\mathrm{uc}}(p).
    \end{equation}

    \item Derive the exact piecewise inner solution $c^*(p)$.

    \item Derive the piecewise value function $\phi(p)$ and identify every
    point at which its derivative changes expression.

    \item Solve the nested problem analytically.

    \item Introduce Lagrange multipliers $\mu_L$ and $\mu_U$ for
    \begin{equation}
    -c\leq0,
    \qquad
    c-1\leq0,
    \end{equation}
    and write the inner KKT conditions:
    \begin{align}
    2(c-p)+0.2c-\mu_L+\mu_U&=0,\\
    \mu_Lc&=0,\\
    \mu_U(c-1)&=0,\\
    \mu_L,\mu_U&\geq0.
    \end{align}

    \item Reformulate the bilevel problem as a single-level mathematical
    program with complementarity constraints.

    \item Enumerate the possible active sets and recover the nested optimum
    from the KKT reformulation.

    \item Solve the direct simultaneous problem
    \begin{equation}
    \min_{0\leq p\leq3,\;0\leq c\leq1}
    (p-2)^2+(c-p)^2+0.1c^2
    \end{equation}
    and compare its solution with the nested result.

    \item Explain why complementarity constraints create numerical
    difficulties for conventional nonlinear programming solvers.
\end{enumerate}


\subsection*{Problem 6: Computational Benchmark of CCD Architectures}

Consider the controlled oscillator
\begin{equation}
\dot{\mathbf{x}}
=
\begin{bmatrix}
0&1\\
-k&-c
\end{bmatrix}
\mathbf{x}
+
\begin{bmatrix}
0\\
1
\end{bmatrix}u,
\end{equation}
with feedback
\begin{equation}
u=-K_p x_1-K_d x_2.
\end{equation}

The design bounds are
\begin{equation}
0.2\leq k\leq5,
\qquad
0.05\leq c\leq3,
\end{equation}
\begin{equation}
0\leq K_p\leq10,
\qquad
0\leq K_d\leq8.
\end{equation}

Define
\begin{equation}
A_{\mathrm{cl}}
=
\begin{bmatrix}
0&1\\
-(k+K_p)&-(c+K_d)
\end{bmatrix},
\end{equation}
and
\begin{equation}
Q=
\begin{bmatrix}
10&0\\
0&1
\end{bmatrix},
\qquad
R=0.05.
\end{equation}

For every stable design, let $P$ satisfy
\begin{equation}
A_{\mathrm{cl}}^TP
+
PA_{\mathrm{cl}}
+
Q
+
\begin{bmatrix}
K_p\\
K_d
\end{bmatrix}
R
\begin{bmatrix}
K_p&K_d
\end{bmatrix}
=
0.
\end{equation}

Use
\begin{equation}
\Sigma_0=
\begin{bmatrix}
1&0\\
0&0.25
\end{bmatrix},
\end{equation}
and define
\begin{equation}
J(k,c,K_p,K_d)
=
\operatorname{tr}(P\Sigma_0)
+
0.02k^2
+
0.02c^2
+
0.005K_p^2
+
0.005K_d^2.
\end{equation}

\begin{enumerate}
    \item Derive the closed-loop stability conditions.

    \item Implement a single-pass sequential architecture:
    \begin{enumerate}
        \item optimize $k$ and $c$ with $K_p=K_d=0$;
        \item freeze the plant and optimize $K_p$ and $K_d$.
    \end{enumerate}

    \item Implement iterative sequential design by alternating exact or
    numerically converged plant and controller subproblems until
    \begin{equation}
    \left\|
    \mathbf{z}^{(j+1)}-\mathbf{z}^{(j)}
    \right\|_2
    <10^{-6}.
    \end{equation}

    \item Implement nested CCD:
    \begin{equation}
    \min_{k,c}
    \left[
    \min_{K_p,K_d}
    J(k,c,K_p,K_d)
    \right].
    \end{equation}

    \item Implement simultaneous CCD:
    \begin{equation}
    \min_{k,c,K_p,K_d}
    J(k,c,K_p,K_d).
    \end{equation}

    \item Use at least ten initial guesses for the nested and simultaneous
    problems. MATLAB \texttt{fmincon}, IPOPT, SNOPT, or an equivalent
    constrained optimizer may be used.

    \item Compare:
    \begin{equation}
    J^*,
    \qquad
    k^*,
    \qquad
    c^*,
    \qquad
    K_p^*,
    \qquad
    K_d^*,
    \end{equation}
    together with function evaluations, derivative evaluations, CPU time,
    and sensitivity to initialization.

    \item Repeat the nested study with inner optimality tolerances
    $10^{-3}$, $10^{-6}$, and $10^{-9}$. Quantify the effect of imperfect
    inner solves on the outer optimum.

    \item Determine whether the best nested and simultaneous solutions agree
    to within
    \begin{equation}
    10^{-5}
    \end{equation}
    in objective value.

    \item Explain any disagreement in terms of local minima, inner-loop
    accuracy, derivative quality, scaling, or solver termination criteria.
\end{enumerate}
